\section*{NeurIPS Paper Checklist}

\begin{enumerate}

\item {\bf Claims}
    \item[] Question: Do the main claims made in the abstract and introduction accurately reflect the paper's contributions and scope?
    \item[] Answer: \answerYes{}
    \item[] Justification: All quantitative claims (92.6\% routing accuracy, 85.3\% ML-only rate, 14.7\% cascade rate) are directly supported by experimental results reported in Section~\ref{sec:experiments}. The system architecture claims are backed by the implementation description in Sections~\ref{sec:architecture}--\ref{sec:personalization}.

\item {\bf Limitations}
    \item[] Question: Does the paper discuss the limitations of the work performed by the authors?
    \item[] Answer: \answerYes{}
    \item[] Justification: Section~\ref{sec:limitations} explicitly discusses limitations including reliance on closed proprietary models, use of synthetic evaluation data, the debug\_simple category weakness (33.3\% accuracy), and lack of cost benchmarking against real-world API pricing.

\item {\bf Theory assumptions and proofs}
    \item[] Question: For each theoretical result, does the paper provide the full set of assumptions and a complete (and correct) proof?
    \item[] Answer: \answerNA{}
    \item[] Justification: This paper presents an empirical systems contribution with no theoretical results or formal proofs. All claims are supported by empirical evaluation.

    \item {\bf Experimental result reproducibility}
    \item[] Question: Does the paper fully disclose all the information needed to reproduce the main experimental results of the paper to the extent that it affects the main claims and/or conclusions of the paper?
    \item[] Answer: \answerYes{}
    \item[] Justification: Section~\ref{sec:experiments} and Appendix~\ref{app:impl} provide all model hyperparameters, feature engineering details, training data sizes, evaluation protocols, and confidence thresholds needed to reproduce results. Code is publicly available at \url{https://github.com/sandeepvijayarao09/dual-llm-system}.

\item {\bf Open access to data and code}
    \item[] Question: Does the paper provide open access to the data and code, with sufficient instructions to faithfully reproduce the main experimental results, as described in supplemental material?
    \item[] Answer: \answerYes{}
    \item[] Justification: The full implementation, 280 seed training examples, 1,500-case evaluation suite, and training scripts are released at \url{https://github.com/sandeepvijayarao09/dual-llm-system} under the MIT License.

\item {\bf Experimental setting/details}
    \item[] Question: Does the paper specify all the training and test details necessary to understand the results?
    \item[] Answer: \answerYes{}
    \item[] Justification: Table~\ref{tab:config} lists all model hyperparameters. Section~\ref{sec:ml_router} describes feature engineering. Section~\ref{sec:experiments} describes the evaluation protocol, category definitions, and confidence thresholds.

\item {\bf Experiment statistical significance}
    \item[] Question: Does the paper report error bars suitably and correctly defined or other appropriate information about the statistical significance of the experiments?
    \item[] Answer: \answerNo{}
    \item[] Justification: The ML router evaluation uses deterministic inference (no randomness), so error bars are not applicable. The evaluation set is fixed and labelled, making repeated runs identical. Future work will evaluate across multiple training seeds to report variance.

\item {\bf Experiments compute resources}
    \item[] Question: For each experiment, does the paper provide sufficient information on the computer resources needed to reproduce the experiments?
    \item[] Answer: \answerYes{}
    \item[] Justification: Section~\ref{sec:experiments} states that all ML router evaluations run on CPU in under 5 seconds. The system uses the OpenAI API for LLM calls; typical API costs are noted in the Limitations section.

\item {\bf Code of ethics}
    \item[] Question: Does the research conducted in the paper conform, in every respect, with the NeurIPS Code of Ethics?
    \item[] Answer: \answerYes{}
    \item[] Justification: This work develops a query routing and personalization system. No human subjects data was collected. Evaluation data is synthetic. No harmful content generation or discriminatory systems are involved.

\item {\bf Broader impacts}
    \item[] Question: Does the paper discuss both potential positive societal impacts and negative societal impacts of the work performed?
    \item[] Answer: \answerYes{}
    \item[] Justification: Section~\ref{sec:limitations} discusses positive impacts (reduced API costs, democratized LLM access) and potential negative impacts (profile data privacy, routing errors for sensitive queries).

\item {\bf Safeguards}
    \item[] Question: Does the paper describe safeguards that have been put in place for responsible release of data or models that have a high risk for misuse?
    \item[] Answer: \answerNA{}
    \item[] Justification: This paper releases a routing system that orchestrates existing commercial LLMs (GPT-4o, GPT-4o-mini). The system does not release a new language model or scraped dataset that poses misuse risk.

\item {\bf Licenses for existing assets}
    \item[] Question: Are the creators or original owners of assets used in the paper properly credited?
    \item[] Answer: \answerYes{}
    \item[] Justification: GPT-4o and GPT-4o-mini (OpenAI API, proprietary license) are cited. scikit-learn (BSD-3 license) is cited. All software dependencies are listed in \texttt{requirements.txt} with version pins.

\item {\bf New assets}
    \item[] Question: Are new assets introduced in the paper well documented?
    \item[] Answer: \answerYes{}
    \item[] Justification: The 280-example seed training dataset and 1,500-example evaluation suite are documented in Section~\ref{sec:experiments} and released with the codebase. The trained model checkpoint (\texttt{router\_v0.joblib}) is included in the repository.

\item {\bf Crowdsourcing and research with human subjects}
    \item[] Question: For crowdsourcing experiments and research with human subjects, does the paper include full instructions given to participants?
    \item[] Answer: \answerNA{}
    \item[] Justification: This paper involves no crowdsourcing or human subjects research. All evaluation data is synthetically generated by the authors.

\item {\bf Institutional review board (IRB) approvals or equivalent for research with human subjects}
    \item[] Question: Does the paper describe potential risks incurred by study participants?
    \item[] Answer: \answerNA{}
    \item[] Justification: No human subjects research was conducted. IRB approval is not applicable.

\item {\bf Declaration of LLM usage}
    \item[] Question: Does the paper describe the usage of LLMs if it is an important, original, or non-standard component of the core methods in this research?
    \item[] Answer: \answerYes{}
    \item[] Justification: LLMs (GPT-4o and GPT-4o-mini via the OpenAI API) are the central components of the proposed system. Their roles as the routing classifier, answer generator, and personalizer are described in detail in Sections~\ref{sec:architecture} and~\ref{sec:personalization}.

\end{enumerate}
